\section{Delete a Node in a BST}
\label{sec:bst-delete}

\problemstatement{Given the root of a BST and a value $X$, delete the
node with value $X$ (if present) and return the root of the resulting
BST, which must still satisfy the BST invariant.}

\begin{patternbox}
\emph{``Delete from a BST''} is the hardest of the three basic operations
because, unlike search or insert, the node being removed might be in the
\textbf{middle} of the tree with two children that both need to stay
connected to the rest of the tree somehow. The keyword to notice: deletion
needs a \textbf{replacement value} that preserves the BST invariant, and
Section~\ref{sec:bst-min-max}'s ``walk to the edge'' primitive is exactly
what finds one.
\end{patternbox}

\subsection*{Understanding the problem carefully}

First, locate the node with value $X$ using the same walk as
Section~\ref{sec:bst-search}. Once found, there are three structurally
different cases:

\begin{itemize}
  \item \textbf{Leaf node (no children).} Simply remove it -- nothing
        else needs to change.
  \item \textbf{One child.} Replace the node with its single child --
        the parent's pointer to this node is redirected to point at the
        child instead, and the BST invariant is automatically preserved,
        since that child's entire subtree already sat correctly on one
        side of the now-removed node.
  \item \textbf{Two children.} This is the hard case: we cannot just
        remove the node, since that would disconnect two subtrees from
        each other. Instead, find a \textbf{replacement value} that can
        legally take this node's place -- the smallest value in the
        right subtree (its inorder successor) or, symmetrically, the
        largest value in the left subtree (its inorder predecessor) both
        work. Copy that replacement value into the current node, then
        recursively delete the replacement's \emph{original} location
        from whichever subtree it came from (which, being an edge value
        of that subtree, will only ever fall into the leaf or one-child
        case, never back into the two-children case).
\end{itemize}

\begin{ideabox}[Two-Children Case: Borrow the Successor, Then Delete It Recursively]
When deleting a node with two children, find its inorder successor (the
minimum of its right subtree, found via Section~\ref{sec:bst-min-max}'s
walk), copy that successor's value into the current node, and then
recursively delete that same value from the right subtree -- which now
contains a duplicate, but only temporarily, and the recursive delete call
removes the original occurrence cleanly.
\end{ideabox}

\subsection*{Why borrowing the successor's value, then deleting it, is always correct}

The inorder successor of a two-children node is, by definition, the
smallest value strictly greater than the node's current value -- which
means it is simultaneously larger than everything in the node's left
subtree (since it's larger than the node itself) and smaller than or
equal to everything else remaining in the right subtree (since it's the
minimum of that subtree). So placing it directly into the current node's
position preserves the BST invariant relative to both subtrees perfectly.
And because the successor is the minimum of the right subtree, it has no
left child (a node with a left child cannot be a subtree's minimum), so
deleting its original location is guaranteed to fall into the simple
leaf-or-one-child case, never recursing into another two-children
scenario at that step.

\subsection*{Algorithm}

\begin{enumerate}
  \item If the current node is null: return null ($X$ was not found;
        nothing to delete).
  \item If $X <$ current node's value: recurse left, reattach the
        result.
  \item If $X >$ current node's value: recurse right, reattach the
        result.
  \item Else (found the node to delete):
  \begin{enumerate}
    \item If it has no left child: return its right child (handles both
          the leaf case and the one-right-child case).
    \item If it has no right child: return its left child.
    \item Otherwise (two children): find $\textit{successor} \gets$ the
          minimum value in the right subtree. Set the current node's
          value to $\textit{successor}$. Recursively delete
          $\textit{successor}$ from the right subtree, and reattach the
          result as the new right child.
  \end{enumerate}
  \item Return the (possibly updated) current node.
\end{enumerate}

\subsection*{Dry run}

\dryrun{Delete $X=5$ from the BST with root $5$, left child $3$, right
child $8$ (with left child $7$ and right child $9$).}

At the root ($5=X$): two children present. Find successor: minimum of
right subtree (rooted at $8$) is reached by walking left: $8 \to 7$ (no
left child) -- successor is $7$. Set root's value to $7$. Recursively
delete $7$ from the right subtree (rooted at $8$): at $8$, $7<8$, go
left to $7$; at $7$ ($7=X$ for this recursive call), no children --
return null, attach as $8$'s new left child.

Resulting tree: root $7$, left child $3$, right child $8$ (now with no
left child, right child $9$).

\subsection*{C++ implementation}

\begin{lstlisting}
#include <bits/stdc++.h>
using namespace std;

struct TreeNode {
    int val;
    TreeNode* left;
    TreeNode* right;
    TreeNode(int v) : val(v), left(nullptr), right(nullptr) {}
};

int findMin(TreeNode* root) {
    while (root->left) root = root->left;
    return root->val;
}

TreeNode* deleteNode(TreeNode* root, int X) {
    if (root == nullptr) return nullptr;

    if (X < root->val) {
        root->left = deleteNode(root->left, X);
    } else if (X > root->val) {
        root->right = deleteNode(root->right, X);
    } else {
        if (root->left == nullptr) return root->right;
        if (root->right == nullptr) return root->left;

        int successor = findMin(root->right);
        root->val = successor;
        root->right = deleteNode(root->right, successor);
    }
    return root;
}

int main() {
    TreeNode* root = new TreeNode(5);
    root->left = new TreeNode(3);
    root->right = new TreeNode(8);
    root->right->left = new TreeNode(7);
    root->right->right = new TreeNode(9);

    root = deleteNode(root, 5);
    cout << "New root: " << root->val << endl;
    return 0;
}
\end{lstlisting}

\begin{complexitybox}
\textbf{Time Complexity.} $O(h)$: at most one $O(h)$ walk to find the
node, plus at most one more $O(h)$ walk to find and delete the successor
in the two-children case.
\textbf{Space Complexity.} $O(h)$ for the recursion stack.
\end{complexitybox}

\begin{takeawaybox}
The hard case of BST deletion is always solved the same way: borrow a
boundary value (successor or predecessor) from the subtree that must stay
connected, copy it up, then delete the boundary value's original location
-- which is guaranteed to be an easy case by construction. This
``borrow an edge value, then recursively clean up'' pattern is worth
remembering as its own primitive, independent of the BST context.
\end{takeawaybox}
